% Independently assembled research artifact; not submitted or accepted.
% Official icml2026.sty and icml2026.bst are used without modification.
% Experimental tables and figures are inputs, not regenerated by this document.
\documentclass{article}
\usepackage{microtype}
\usepackage{graphicx}
\usepackage{booktabs}
\usepackage{amsmath,amssymb,amsthm,mathtools}
\usepackage{hyperref}
\newcommand{\theHalgorithm}{\arabic{algorithm}}
\usepackage{icml2026}
\usepackage[T1]{fontenc}
\newtheorem{proposition}{Proposition}
\newcommand{\R}{\mathbb{R}}
\newcommand{\E}{\mathbb{E}}
\newcommand{\tr}{\operatorname{tr}}
\newcommand{\Cov}{\operatorname{Cov}}
\newcommand{\rank}{\operatorname{rank}}
\newcommand{\row}{\operatorname{row}}
\newcommand{\artifact}[1]{\path{#1}}
\icmltitlerunning{Measure the Difficult, Predict the Rest}
\begin{document}
\twocolumn[
\icmltitle{Measure the Difficult, Predict the Rest:\\Risk-Aware Benchmark Selection and Its Limits}
\begin{icmlauthorlist}
\icmlauthor{Anonymous Artifact Author}{artifact}
\end{icmlauthorlist}
\icmlaffiliation{artifact}{Independent anonymous research artifact}
\icmlcorrespondingauthor{Anonymous Authors}{}
\icmlkeywords{benchmark selection, score prediction, experimental design, conditional Gaussian models}
]
\printAffiliationsAndNotice{This manuscript is a research artifact and has not been submitted or accepted.}

\begin{abstract}
Can a few complete benchmarks predict a model's scores on the rest? We study an independently assembled conditional-Gaussian selector whose squared-risk objective rewards both information about unmeasured tasks and removal of difficult tasks from evaluation. On an availability-selected, 47-task MTEB embedding cohort, measuring eight tasks reduces mean per-model unrun-score RMSE from 4.42 to 3.45 percentage points relative to random selection with the same predictor; a secondary six-task generative leaderboard shows 7.20 versus 4.36 points at two tasks. Fixed hidden targets reveal genuine but smaller transfer, while validation-supervised greedy selection often performs better and additional support shrinkage has no supported benefit. Repeated-split, publisher-clustered descriptive intervals, missingness stresses, duplicate audits, and exploratory target-set checks delimit these findings rather than certify a universal sufficient budget. The method uses standard statistical components, and neither its Gaussian variance nor its empirical threshold crossings provide calibrated prediction guarantees.
\end{abstract}

\section{Introduction}
Running every benchmark for every model is expensive. Historical model--task score matrices suggest a practical alternative: measure a few complete tasks and predict the remaining aggregate scores. This is a different budget from choosing a few individual examples within every benchmark. It also poses a subtle evaluation problem: a policy that measures the most difficult-to-predict tasks no longer needs to predict them. Lower error on the remaining tasks can therefore reflect both useful information and a change in what is scored.

We make this distinction central. A conditional-Gaussian model supplies a simple, train-only ordering of benchmark queries. Its exact one-step squared-risk reduction decomposes into the variance of the queried task, which disappears from the target set, and covariance-mediated transfer to still-hidden tasks. An independent-coordinate example shows that the first term alone can produce an apparent improvement without learning anything about the remaining scores. A fixed, unqueryable target set removes this mechanism and directly tests transfer.

Our primary evidence concerns English \emph{embedding} models on MTEB; a six-task generative-model leaderboard is a separate, narrower secondary setting. Both are retrospective, availability-selected complete-score evaluation cohorts. They do not identify accuracy on naturally unreported scores, and publisher or name-based grouping does not guarantee unseen model genealogies. These limitations matter even when test scores are not used for fitting.

The contribution is an empirical and mathematical audit, not a claim of first invention or state-of-the-art performance. We (i) specify a reproducible standard covariance-design and prediction pipeline; (ii) separate moving-target error from fixed-target transfer, including stronger selector--predictor controls; and (iii) answer ``how many tasks are enough?'' only for explicit tolerances, cohorts, and tested budgets. The positive result is that historical design improves on random selection with the same predictor. The negative results are equally consequential: extra support shrinkage is unsupported, supervised greedy is often stronger, and small-budget success depends heavily on the targets and historical support.

\begin{figure*}[t]
\centering
\includegraphics[width=\textwidth]{artifacts/figures/budget_curves.pdf}
\caption{Budget curves on repeated random/native splits: primary MTEB embeddings (left), secondary OpenLLM generative tasks (right). Bands for risk design and random covariance are conditional 95\% bootstrap intervals; dotted/dashed lines mark illustrative 3/5 pp tolerances. Only unrun tasks are scored, so larger $k$ changes target composition. Low-rank risk design is exploratory. Source: \artifact{artifacts/figures/budget_curves.pdf}.}
\label{fig:budget}
\end{figure*}

\section{Related Work and Scope}
\label{sec:related}
\paragraph{The unit of selection.}
MTEB evaluates heterogeneous embedding tasks with different metrics and task-specific strengths \citep{muennighoff2022mteb}. A column here is one specified aggregate score, not an interchangeable evaluation item. \citet{polo2024tinybenchmarks} select examples and use item-response models to reconstruct scenario scores. \citet{kipnis2024metabench} select 739 items across six source benchmarks and use their inferred abilities jointly. These results motivate economical evaluation but do not establish a sufficient number of \emph{entire} benchmarks to omit. Our aggregate matrices do not contain the item responses required to reproduce those IRT procedures.

\paragraph{Prediction and agreement.}
\citet{ruan2024observationalscaling} predict downstream performance from low-dimensional observed benchmark capabilities; their experimental-design budget selects historical model rows, not query columns for a new model. \citet{polo2024sloth} predict performance across families using resources and smaller same-family anchors, a different information setting. These works motivate low-rank controls but are not reproduced here. Benchmark agreement depends on the model set, granularity, and metric \citep{perlitz2024benchbench}; high ranking correlation alone does not certify accurate missing-score prediction.

\paragraph{Restricted positioning.}
Only six supplied papers, in their archived v1 versions, were available for related-work reading. No broader priority check, citation traversal, or search for the excluded exact-problem work or work citing it was conducted. The method is an independently assembled application of standard conditioning and greedy design. We make no exhaustive novelty, universal superiority, or priority claim. Appendix~\ref{app:access} documents this access constraint.

\section{Risk-Aware Selection and Prediction}
\label{sec:method}
\subsection{Problem and Predictor}
Let $X\in\R^{M\times n}$ contain historical model scores, possibly missing. A new model has score vector $x$. We select $S\subset V=\{1,\ldots,n\}$, $|S|=k<n$, observe only $x_S$, and predict $x_U$ for $U=V\setminus S$. The operational budget counts complete columns at equal cost; it does not measure tokens, dollars, or wall time.

Fit training-only means $\hat\mu$, scales, and a positive-definite covariance $\hat C$. The pre-clipping conditional predictor is
\begin{equation}
\hat x_U=\hat\mu_U+\hat C_{US}\hat C_{SS}^{-1}(x_S-\hat\mu_S).
\label{eq:predict}
\end{equation}
Predictions are clipped to theoretical metric bounds, not empirical extrema. Correlation-based MTEB tasks permit $[-1,1]$; the other included scores use $[0,1]$. The fitting recipe standardizes observed entries, estimates pairwise second moments, optionally shrinks off-diagonal entries by $n_{ab}/(n_{ab}+\lambda)$, repairs eigenvalues, renormalizes the diagonal, and mixes with $\alpha I$. We call this pipeline support-regularized risk design (SRRD), even when validation chooses $\lambda=0$. Pairwise products use per-task training means, not pair-specific recentering; they are an estimation heuristic, not an identified missing-data likelihood. Exact details are in Appendix~\ref{app:fit}.

\subsection{Why Measuring Hard Tasks Can Help}
For known $\mu$ and $C\succ0$ with $x\sim\mathcal N(\mu,C)$, let
\begin{equation}
C^S=C-C_{:S}C_{SS}^{-1}C_{S:}.
\label{eq:conditional}
\end{equation}
The conditional mean minimizes expected summed squared error and has unrun risk $R(S)=\tr(C^S_{UU})$. Querying $j\in U$ updates covariance by
\begin{equation}
C^{S\cup\{j\}}=C^S-\frac{C^S_{:j}C^S_{j:}}{C^S_{jj}}.
\label{eq:update}
\end{equation}
Consequently its exact one-step gain is
\begin{equation}
\Delta_U(j\mid S)=\underbrace{C^S_{jj}}_{\text{target removal}}+
\underbrace{\sum_{t\in U\setminus\{j\}}\frac{(C^S_{tj})^2}{C^S_{jj}}}_{\text{transfer}}.
\label{eq:gain}
\end{equation}
SRRD greedily maximizes this gain using $\hat C$. Variance-only selection retains the first term; transfer-only selection retains the second. Paths depend on the fitted covariance, not the new model's observed score values. They are sequential computation of a precomputed ordering, not model-specific adaptive querying.

For a fixed hidden set $T$ and query pool $V\setminus T$, the corresponding objective is $R_T(S)=\tr(C^S_{TT})$, with gain
\begin{equation}
\Delta_T(j\mid S)=\sum_{t\in T}\frac{(C^S_{tj})^2}{C^S_{jj}},\qquad j\notin S\cup T.
\label{eq:fixedgain}
\end{equation}
There is no removal term. Target-risk design uses this objective; candidate-restricted unrun-risk design uses Equation~\eqref{eq:gain} on the same query pool, and both are scored only on $T$.

\paragraph{A sharp limit.}
With $C=\operatorname{diag}(9,1)$, measuring coordinate 1 reduces unrun summed risk from 10 to 1 but provides zero information about fixed target $T=\{2\}$. Conversely, in a noiseless known-loading model $x=\mu+Bz$, recovery of targets $A$ is possible exactly when $\row(B_A)\subseteq\row(B_S)$. Small $k$ thus requires informative shared structure, not just a favorable moving-target curve. Appendix~\ref{app:proofs} gives full derivations and the low-rank converse.

\paragraph{What is not guaranteed.}
The identities concern known Gaussian parameters, exact arithmetic, and \emph{expected squared error}. Our metric averages per-model RMSE; this is not the square root of expected MSE. The exact selector correspondence requires the implementation's $10^{-14}$ denominator guard to remain inactive. The Gaussian mean interpretation is pre-clipping only. Estimated posterior variances are not calibrated prediction intervals, and no global greedy-optimality, submodularity, or approximation guarantee is claimed.

\begin{figure*}[t]
\centering
\includegraphics[width=\textwidth]{artifacts/figures/fixed_target_curves.pdf}
\caption{Fixed-target transfer on the original pre-outcome seed-701 audit. Target-risk and candidate-restricted unrun-risk design use the same fitted covariance; supervised greedy uses ridge. Unlike Figure~\ref{fig:budget}, targets cannot be removed. The OpenLLM pair is IFEval/MMLU-PRO, not an outcome-optimized choice. Source: \artifact{artifacts/figures/fixed_target_curves.pdf}.}
\label{fig:fixed}
\end{figure*}

\section{Evaluation Protocol}
\label{sec:protocol}
\subsection{Data and Quality Control}
\paragraph{Primary: MTEB embeddings.}
We use results commit \texttt{6dabbcc68c28ea5f180c\allowbreak e7df79347b3ca9cc6383}. From a canonical 56-task English list, extraction retains exact English-only \texttt{main\_score} records on test (MSMARCO uses dev), averages eligible subsets, and selects one revision per model directory by valid coverage with a lexical tie-break. There is no metric fallback or cross-revision stitching. Explicit nontext and cross-encoder revisions are excluded. Among 217 models with at least 20/56 scores, the 47 columns observed for at least 160 models define the primary universe: 131 complete rows are eligible for validation/test and 86 partial rows augment historical training only. The nine excluded tasks receive no generalization claim. Only 56 models complete all 56 tasks; the optional all-56 predictive experiment was not run.

\paragraph{Secondary: six generative tasks.}
The Open LLM Leaderboard (OpenLLM) snapshot is commit \texttt{9c09a7cae43334062a82\allowbreak cb164f2ef255013dafa2}. We use raw $[0,1]$ scores for IFEval, BBH, MATH Lvl 5, GPQA, MUSR, and MMLU-PRO, not normalized scores or leaderboard averages. Of 4,576 source rows, 4,575 pass complete/unflagged/range filtering; deterministic exact-fullname deduplication leaves 4,496 rows. No identities are joined across leaderboards.

Task universes and evaluation eligibility use full-snapshot \emph{availability masks}; this is mask-transductive cohort construction, not score-value-driven task discovery or evidence of random missingness. Snapshot heterogeneity and training-data contamination remain possible. Cell provenance, task coverage, range checks, identity handling, and exact-score duplicates are audited. In particular, 100 OpenLLM rows belong to repeated complete six-score vectors; primary random splits have 10--16 train/test-shared distinct vectors, and publisher splits have 2--4 except one with zero. We do \emph{not} claim the primary evaluation is leakage-free. MTEB has no repeated complete 47-score vectors in this audit. Appendix~\ref{app:duplicates} reports the completed exclusion sensitivity.

\subsection{Splits, Comparisons, and Stages}
Seeds $11,23,37,53,71$ define five repeated random approximately 60/20/20 train/validation/test splits, grouping canonical identities and embedding-dimension variants. These splits overlap and are not independent replications. Publisher-disjoint and mechanical name-family groupings stress generalization; neither verifies backbone or genealogical separation. Partial MTEB rows sharing held-out identities/groups are excluded from training. Means, scales, imputation, covariance, latent factors, and regression coefficients use training entries only. Predictors receive only selected test coordinates; hidden scores enter evaluation only.

We compare task means, standardized per-model quality offsets, random queries with direct ridge, missing-data low-rank, or the same covariance predictor, and validation-forward-greedy ridge. The latter shares validation data between predictor tuning and query selection: it is a strong but overfit-prone supervised baseline, not nested cross-validation. Ten deterministic random permutations are averaged \emph{within each model's loss}. Mean-on-risk-targets, no-support, minimally repaired pairwise, variance-only, and transfer-only controls diagnose mechanisms. Crossed ridge/covariance controls and selected low-rank design were added after the core as exploratory frozen-parameter checks, without new hyperparameter search.

The pre-outcome plan had four stages: data/baseline integrity, bounded validation tuning, core method comparison, and stress/ablation. Success gates were valid held-out-model evaluation and complete reporting, not positive gains. Hypotheses concerned design versus random with the same predictor (H1), added support shrinkage (H2), cohort-dependent sufficient budgets (H3), and removal versus transfer (H4). Binding amendments supersede the original plan's population-testing language. A first-seed pilot led only to logging/storage optimization; stale pilot outputs are excluded from final aggregation.

\subsection{Metrics, Uncertainty, and Budgets}
For model $i$, split $s$, method $a$, and query repetition $r$, the primary loss is
\begin{equation}
\ell_{siar}=100\sqrt{\frac{1}{|T_{sar}|}\sum_{t\in T_{sar}}
(\hat x_{siart}-x_{sit})^2},
\label{eq:metric}
\end{equation}
where $T$ is either the unrun set or a fixed hidden target set. The estimand is the equal average of five split-level means after random repetitions are averaged within each model. Secondary metrics are MAE, training-scale-standardized RMSE, and the average of split-level 90th percentiles of model RMSE. Percentage points (pp) are score units, not relative percentages or equal downstream utility.

We jointly resample publisher clusters across every repeated model/split occurrence, recomputing the \emph{exact equal-split} statistic in 3,000 bootstrap draws. Paired contrasts share resample weights. Intervals condition on fixed splits and fitted predictors: they exclude refitting/training uncertainty and are not prediction intervals. Split SD is reported separately as descriptive split/fit sensitivity. All comparisons are descriptive, with no confirmatory population $p$-values or multiple-testing significance claims; intervals across methods, budgets, and targets are not simultaneous bounds.

Budgets are $\{1,2,4,8,12,16,24,32\}$ of 47 MTEB tasks and $\{1,2,3,4,5\}$ of six OpenLLM tasks. Prechosen illustrative tolerances are 3 and 5 pp. We report the first \emph{tested} mean crossing, not an interpolated, universal, or confidence-certified minimum. Each moving-target budget changes the scored tasks. Fixed targets, drawn once before outcomes with seed 701, contain 12 MTEB tasks and the OpenLLM pair IFEval/MMLU-PRO; the latter permits at most four candidate queries.

\section{Results}
\label{sec:results}
\begin{table*}[t]
\caption{Unrun-task mean per-model RMSE (pp), with conditional 95\% publisher-cluster bootstrap intervals in brackets. Lower is better. Random and publisher denote held-out-model split rules, not query policies. The last two rows and low-rank risk design are post-core exploratory frozen-parameter controls. These ten methods vary both selection and prediction; the random-covariance comparison isolates selection policy but still changes unrun target composition.}
\label{tab:main}
\centering
\input{artifacts/tables/main_results.tex}
\end{table*}

\subsection{Train-Only Design Helps Against Random Selection}
At the primary MTEB budget $k=8$, SRRD achieves 3.4466 pp versus 4.4206 for random queries with the identical covariance predictor (Table~\ref{tab:main}). The paired difference is $-0.9740$ pp with conditional 95\% interval $[-1.1531,-0.8080]$; split SDs are 0.3298 and 0.3579 pp, respectively. In the secondary OpenLLM setting at $k=2$, the corresponding scores are 4.3552 versus 7.1999 pp, a difference of $-2.8447$ pp $[-3.2344,-2.4833]$, with split SDs 0.0529 and 0.5570 pp. These support H1 descriptively, not as population significance tests.

The result is not dominance over stronger alternatives. Validation-greedy ridge obtains 3.6253 pp on MTEB and 4.3062 on OpenLLM. In the exploratory crossed controls, risk-selected queries with ridge achieve 3.4264 pp on MTEB, whereas ridge-selected queries with covariance achieve 3.7118. Selected low-rank design obtains 3.7194, far better than random low-rank at 4.4544. Thus a weak random low-rank baseline alone would not identify a covariance-specific advantage. Publisher-heldout MTEB scores are higher (SRRD 4.0656 versus random covariance 4.7613), while low-rank design and risk queries with ridge are slightly lower than SRRD. There is no uniform ordering of all selectors and predictors.

\subsection{Removal Matters, but Fixed-Target Transfer Remains}
\begin{table*}[t]
\caption{Original seed-701 fixed-target audit on random/native splits. Entries are mean model RMSE (pp) [conditional 95\% interval]. The 12 MTEB targets or two OpenLLM targets remain hidden at every budget, and every method uses the same disjoint candidate pool. Low-rank risk design is exploratory. Scores are not directly subtractable from Table~\ref{tab:main} to estimate a causal removal fraction.}
\label{tab:fixed}
\centering
\input{artifacts/tables/fixed_results.tex}
\end{table*}

Table~\ref{tab:fixed} isolates transfer. Target-risk design achieves 3.7322 versus random covariance 4.1875 pp on MTEB at $k=8$, a paired difference of $-0.4553$ pp $[-0.6794,-0.2551]$. On the OpenLLM hidden pair at $k=2$, it achieves 9.0429 versus 10.0173 pp, a difference of $-0.9744$ pp $[-1.1485,-0.8108]$. Transfer is therefore present, but the absolute fixed-target error is substantial, and supervised greedy is lower on both audits (3.2666 and 8.9661 pp). Target-aware selection is not consistently better than candidate-restricted unrun-risk selection.

A no-information control also shows a target-composition effect: retaining training-mean predictions but scoring only the risk policy's remaining targets reduces MTEB error from 9.8756 (random targets) to 8.3127 pp, and OpenLLM from 11.4027 to 8.5869 pp. No measured values enter these predictions; these contrasts concern target composition, not a decomposition of the full predictor's gain.

In all 15 native OpenLLM split/seed configurations, the main $k=2$ risk selector measures IFEval and MMLU-PRO---exactly the pair that seed 701 had independently fixed as hidden targets. Their removal makes the main-versus-fixed contrast especially revealing. The pair was one random draw chosen \emph{before} outcomes, not optimized afterward. This does not establish how much of the main gain is causally due to removal: the two audits score different targets under different allowed queries.

Additional target checks were frozen after the core outcomes and are explicitly exploratory. Averaging five MTEB target draws (701--705) gives target-risk/random covariance RMSE 3.3863/4.0567 pp at $k=8$, a difference of $-0.6704$ $[-0.8544,-0.5061]$. Averaging all 15 OpenLLM hidden pairs gives 5.6896/6.3576 pp at $k=2$, a difference of $-0.6680$ $[-0.7356,-0.6017]$. These averages do not turn overlapping target sets into independent datasets. Appendix~\ref{app:sensitivity} reports heterogeneity and the random-permutation distinction from the original audit.

\subsection{Shrinkage and Missingness: Negative Results}
\begin{table*}[t]
\caption{Random/native ablations: mean RMSE $\pm$ split SD (pp), \emph{not} confidence intervals. No-support covariance is tuned over the same diagonal-shrinkage grid with $\lambda=0$. Minimal PSD repair fixes $\lambda=\alpha=0$ but retains the positive eigenvalue floor. Variance and transfer-only rules retain the corresponding terms of Equation~\eqref{eq:gain}.}
\label{tab:ablation}
\centering
\input{artifacts/tables/ablations.tex}
\end{table*}

Additional support shrinkage has no supported benefit (H2). Validation selects $\lambda=0$ in 58/65 final jobs, including every random/native and name-family/native job; SRRD consequently equals the tuned no-support alternative in Table~\ref{tab:ablation}. Native MTEB pairwise estimates are indefinite before repair, but repair necessity is not evidence that support-count shrinkage improves prediction. Transfer-only design is worse on the moving-target endpoint, particularly OpenLLM (7.9956 pp versus SRRD 4.3552); this is compatible with the value of target removal, not proof that transfer is absent.

Historical augmentation is not uniformly helpful. Complete-only MTEB training yields 3.1433 pp at $k=8$, \emph{better} than native training with partial rows (3.4466). This changes both sample size and training population, so it is not a clean isolated missingness effect. Refitting/retuning with 50\% or 25\% random retention of observed training entries gives 4.6442 and 5.3492 pp; deleting each task's upper half yields 4.8816. OpenLLM upper-tail deletion increases error from 4.3552 to 6.3736 pp. These hypothetical masks do not identify or correct natural missing-not-at-random reporting. Figure~\ref{fig:stress} in the appendix shows all five-seed conditions.

\subsection{How Many Benchmarks Are Enough?}
Figure~\ref{fig:budget} answers only the frozen descriptive question. MTEB SRRD first meets 5 pp at $k=4/47$ and 3 pp at $k=16/47$. Its $k=12$ mean is 3.0046 pp, \emph{above} 3; rounding it to 3.00 must not move the crossing to 12. OpenLLM first meets 5 pp at $2/6$ and 3 pp at $3/6$. These are not confidence-certified minimum budgets, and decreasing error can partly reflect removing difficult targets.

Under the frozen 12-target MTEB audit, target-risk SRRD never reaches 3 pp over tested budgets through 32; supervised greedy first crosses at 16 (2.9991 pp). The frozen OpenLLM pair never reaches even 5 pp, including after all four candidates are measured (covariance: 8.7260 pp). The full cohort/method threshold ledger is \artifact{artifacts/analysis/budget_thresholds.csv}. There is no universal sufficient $k$.

\section{Discussion and Limitations}
\label{sec:discussion}
The strongest evidence is for a policy-level design benefit relative to random selection under the same predictor and snapshot distribution. Fixed-target audits demonstrate actual transfer, but also show why unrun-task curves should not be described as pure information gains. The Gaussian decomposition explains a possible mechanism exactly; it neither quantifies that mechanism's empirical share nor validates Gaussian modeling of heterogeneous scores.

Several issues restrict application. Complete-case test cohorts and snapshot-defined task coverage may select unusually well-evaluated models; unreported cells are not validated targets. Publisher identities and name-family keywords are incomplete lineage controls. Exact-score duplicates cross the primary OpenLLM splits; removing all repeated-vector rows is reassuring only within a changed, outcome-defined cohort. Leaderboard measurement noise, evaluator versions, contamination, and related submissions can all create dependence not captured by publisher clusters. There is no temporal or prospective evaluation, item-level IRT reproduction, uncertainty calibration, or verified extrapolation to new benchmark types. The low-rank imputation often does not converge within its 25-iteration budget, so it is a bounded-compute alternative, not an optimized low-rank oracle.

Equal task counts also conceal heterogeneous execution costs: one retrieval benchmark can be far more expensive than one similarity task. Moreover, a pp error on distinct metrics is not the same utility loss. Practical deployment needs cost-aware objectives, reliable lineage metadata, prospective validation, and independently calibrated residual uncertainty. These are future requirements, not results of the present artifact. The full-56 MTEB optional run was not completed, and no inference is made about the nine excluded tasks. The restricted related-work access prevents a complete priority comparison.

\section{Conclusion}
A standard conditional-covariance design can use historical scores to outperform random benchmark selection. Its objective rewards two effects that must be separated: measuring difficult targets and transferring information to genuinely hidden targets. The fixed-target and stronger-control results support useful but limited transfer, reject a supported benefit for extra support shrinkage, and leave supervised greedy competitive or better. Sufficient budgets are descriptive properties of a target set, predictor, cohort, and tolerance---not universal constants or calibrated guarantees.

\section*{Impact Statement}
This work studies reducing aggregate benchmark evaluation effort, not improving a model's underlying capabilities or certifying deployment safety. Replacing measurements with predictions can obscure specialist failures, amplify historical reporting bias, and create misleading leaderboard rankings, particularly for poorly represented models or capabilities. Predicted scores should be labeled as predictions, kept separate from measured scores, and not used as safety or fairness certification. No foundation-model training, paid inference API, or human-subject study was performed; only score predictors were fitted. Public benchmark scores may inherit contamination and consent/licensing issues from their sources; MTEB results are CC0-1.0, while the secondary leaderboard snapshot does not declare a dataset license in the inspected metadata. Public accessibility is not permission to assume unrestricted redistribution. Resource savings here concern a hypothetical future evaluation policy; they are not measured end-to-end benchmark-execution savings.

\bibliography{references}
\bibliographystyle{icml2026}

\clearpage
\appendix
\onecolumn
\section{Full Derivations and Exact Scope}
\label{app:proofs}
These derivations reproduce the audited statements in \artifact{artifacts/theory/theory.md}; the audit is \artifact{artifacts/theory/audit.md}. The local audit used sequential review lenses by the same assistant, not independent reviewers or a formal theorem prover. The statements below apply to known population parameters, not fitted leaderboard correctness.

\subsection{Gaussian Conditioning and Bayes Squared Risk}
Let $x\sim\mathcal N(\mu,C)$ with symmetric $C\succ0$, $V=\{1,\ldots,n\}$, and fixed distinct queried coordinates $S$. For $S=\varnothing$, set $C^S=C$ and $m^S=\mu$. Otherwise define
\begin{align}
m^S&=\mu+C_{:S}C_{SS}^{-1}(x_S-\mu_S),\\
C^S&=C-C_{:S}C_{SS}^{-1}C_{S:}.
\end{align}
\begin{proposition}[Conditional law and risk]
The conditional law of $x$ given $x_S$ has mean $m^S$ and covariance $C^S$, with zero queried rows/columns and positive-definite unqueried block $C^S_{UU}$ when $U\ne\varnothing$. The Bayes conditional and unconditional expected summed squared losses on $U$ both equal $\tr(C^S_{UU})$.
\end{proposition}
\begin{proof}
Every principal $C_{SS}$ is positive definite: extend a nonzero vector on $S$ by zero coordinates and apply $C\succ0$. Put $Y=x-\mu$ and $e=Y-C_{:S}C_{SS}^{-1}Y_S$. Covariance expansion yields $\Cov(e)=C^S$ and $\Cov(e,Y_S)=0$. The pair is jointly Gaussian; its characteristic function factors with zero cross-covariance, proving independence even though the embedded $e$ is singular. Thus $x=m^S+e$ proves the conditional law, and $e_S=0$. For nonzero $w\in\R^{|U|}$, let $v_U=w$ and $v_S=-C_{SS}^{-1}C_{SU}w$. Then $w^TC^S_{UU}w=v^TCv>0$, giving Schur positivity. For any finite-risk measurable predictor $g(x_S)$,
\begin{equation}
\E[\|x_U-g\|^2\mid x_S]=\tr(C^S_{UU})+\|m^S_U-g\|^2,
\label{eq:orthogonal}
\end{equation}
because the conditional cross term vanishes. The conditional mean attains the minimum, uniquely up to almost-sure equality; an infinite-risk estimator cannot improve it. Taking expectations gives the unconditional minimum since $C^S$ does not depend on observed values. Empty $U$ gives zero trace and zero vector loss. These are conditional-version, almost-sure statements, not conditioning on arbitrary null events.
\end{proof}

\subsection{One-Step Gains and Changing Denominators}
\begin{proposition}[Removal and transfer]
For $j\in U$, $v=C^S_{jj}>0$, querying $x_j$ gives Equation~\eqref{eq:update}. The reduction of unnormalized unrun risk is Equation~\eqref{eq:gain}; for fixed $T$ disjoint from all queries it is Equation~\eqref{eq:fixedgain}.
\end{proposition}
\begin{proof}
Set $e'=e-C^S_{:j}e_j/v$. Expansion gives $\Cov(e')=C^S-C^S_{:j}C^S_{j:}/v$, $\Cov(e',e_j)=0$, and $\Cov(e',Y_S)=0$. Joint Gaussianity implies independence of $e'$ from $(Y_S,e_j)$. Given $x_S$, querying $x_j$ reveals $e_j=x_j-m^S_j$, so the new mean is $m^S+C^S_{:j}(x_j-m^S_j)/v$ and the new covariance is the claimed update. For $W=U\setminus\{j\}$,
\begin{equation}
R(S)=v+\sum_{t\in W}C^S_{tt},\qquad
R(S\cup\{j\})=\sum_{t\in W}\left(C^S_{tt}-\frac{(C^S_{tj})^2}{v}\right).
\end{equation}
Subtracting proves the removal-plus-transfer identity. Subtracting traces on the same fixed $T$ gives the fixed-target identity without removal. Each gain is nonnegative. Fixed-target gain vanishes exactly when $C^S_{Tj}=0$, which under the Gaussian law is conditional independence of the query and target vector. Covariances are value-independent, so these also equal gains averaged over the prospective query value.
\end{proof}
For nonempty $U$, expected MSE equals $R(S)/|U|$. At a single step every available candidate leaves the same denominator $|U|-1$, when positive, so next-step expected MSE has the same ordering as the trace objective. Across budgets the denominator changes, and normalized error need not be monotone. The empirical population analogue is instead
\begin{equation}
100\E\sqrt{\frac{\|x_U-m^S_U\|^2}{|U|}}
\ \le\ 100\sqrt{\frac{R(S)}{|U|}},
\end{equation}
by Jensen's inequality. Equality and identical ordering of mean RMSE do not follow. At $k=n$ RMSE on unrun tasks is undefined, not zero, and the scorer rejects empty targets.

For $C=\operatorname{diag}(9,1)$, $S=\varnothing$, querying coordinate 1 reduces summed risk $10\to1$ solely by removal. For fixed $T=\{2\}$ its gain is zero and $x_2\mid x_1\sim\mathcal N(0,1)$. More generally a diagonal positive-definite covariance has positive moving-target removal gains and zero transfer to any disjoint fixed targets.

\subsection{Noiseless Low-Rank Recovery: Sufficiency and Necessity}
\begin{proposition}[Row-space recovery]
Replace positive definiteness by the noiseless model $x=\mu+Bz$, with known $B\in\R^{n\times r}$ and $z\sim\mathcal N(0,I_r)$. Let $L=\row(B_S)$ and let $P$ be its orthogonal projector. Then
\begin{align}
\E[x\mid x_S]&=\mu+BB_S^+(x_S-\mu_S),\\
\Cov(x\mid x_S)&=B(I_r-P)B^T.
\end{align}
Exact recovery of coordinates $A$ by a measurable function of $x_S$, almost surely, is possible if and only if
\begin{equation}
B_A(I_r-P)=0\quad\Longleftrightarrow\quad
\row(B_A)\subseteq\row(B_S)\quad\Longleftrightarrow\quad
\ker B_S\subseteq\ker B_A.
\end{equation}
If $\rank(B_S)=r$, every coordinate is recovered for every $z$ by
$z=(B_S^TB_S)^{-1}B_S^T(x_S-\mu_S)$. When $B$ has full column rank, this rank condition is also necessary for all-coordinate recovery. For arbitrary $B$, the precise condition is $\rank(B_S)=\rank(B)$.
\end{proposition}
\begin{proof}
Decompose $z=Pz+(I_r-P)z$. The two Gaussian components have zero cross-covariance and are independent. Since $B_S(I_r-P)=0$, observations reveal $B_SPz$. The restriction of $B_S$ to $L$ is injective because $L\cap\ker B_S=L\cap L^\perp=\{0\}$. Thus the observations determine $Pz=B_S^+(x_S-\mu_S)$ on their support; the other component retains mean zero and covariance $I_r-P$. This proves the conditional formulas. The minimum target squared risk, by Equation~\eqref{eq:orthogonal}, is
\begin{equation}
\tr[B_A(I_r-P)B_A^T]=\|B_A(I_r-P)\|_F^2.
\end{equation}
It vanishes exactly for the stated row-space inclusion; taking orthogonal complements gives the kernel inclusion. A nonzero conditional risk rules out arbitrary measurable almost-sure recovery, not merely linear reconstruction. The inclusion also yields deterministic recovery for every latent vector. If it fails, there is $h\in\ker B_S$ with $B_Ah\ne0$, so $z$ and $z+h$ agree on observations but differ on targets. Full column rank of $B_S$ makes $P=I_r$ and supplies the left inverse. If $B$ has full column rank, $\row(B)=\R^r$ so full recovery requires $\row(B_S)=\R^r$. Otherwise $\row(B_S)\subseteq\row(B)$ always, and recovery requires equality, equivalently equal ranks.
\end{proof}
Empty $S$ gives $P=0$; only deterministic targets are recoverable. The $r=0$ case is deterministic throughout. For rows $(1,0),(2,0)$ with latent dimension $r=2$, a single queried row recovers all coordinates despite rank $1<r$, illustrating why the converse needs full column rank of $B$. Singular $BB^T$ is not covered by importing the positive-definite inverses from the preceding propositions. The fitted low-rank baseline has strictly positive residual diagonal noise; it is not this noiseless model.

\subsection{Implementation Qualifications}
In \artifact{src/methods.py}, \texttt{Conditional.path} uses $\max(C^S_{jj},10^{-14})$ in denominators, symmetrizes updates, and zeros queried rows/columns. In exact arithmetic and with every candidate guard inactive, induction gives the conditional-covariance invariant and the stated one-step scores. Positive definiteness alone does not guarantee an inactive guard: for $C=\operatorname{diag}(\epsilon,1)$ with $0<\epsilon<10^{-14}$, the guarded query-1 score is $\epsilon^2/10^{-14}$ rather than $\epsilon$. The prediction API clips nonempty-selection outputs; only the pre-clipping expression is the Gaussian conditional mean. Its empty-selection branch returns $\mu$ directly. Callers must explicitly exclude fixed targets from candidates: the path API's target argument does not itself enforce this. The audited experiment callers do supply disjoint pools. No empirical Gaussianity, covariance-estimator consistency under natural missingness, variance calibration, or implementation-wide exact optimality theorem follows.

\subsection{Simulated Population Checks, Not Leaderboard Results}
The actual outputs \artifact{artifacts/theory/theory_mc.csv} and \artifact{artifacts/theory/theory_checks.json} are explicitly synthetic. With 100,000 Gaussian draws per seed $11,23,37$, the theoretical summed risks before/after one query are 3.514071/2.359575. The seed-11 empirical risks are 3.504004/2.360191 with Monte Carlo standard errors 0.007974/0.005814; the other seeds are archived in the CSV. The exact analytic gain is 1.154496, comprising removal 1.056338 and transfer 0.098159. A rank-three noiseless simulated recovery check at $k=3$ gives maximum absolute reconstruction error $1.0658\times10^{-14}$. These verify population algebra numerically, not real-data predictive accuracy or Gaussian fit. Floating-point residual traces near $10^{-14}$ are numerical cancellation, not positive irreducible population risk in the recovered noiseless case.

\section{Data Extraction, Identity Rules, and Splits}
\label{app:data}
\subsection{Pinned Inputs and Availability Selection}
The fixed source endpoints are the MTEB results repository at the commit in Section~\ref{sec:protocol} and the OpenLLM contents dataset at its stated commit. Acquisition is implemented in \artifact{src/download_data.py}; checksums and download provenance are archived under \artifact{artifacts/provenance/}. The MTEB recursive tree response was truncated, so extraction used the full pinned codeload archive instead. The log reports a 395,746,943-byte archive, 88,349 extracted JSON files, and 547 actual model directories, rather than the task brief's approximate 548. These are acquisition counts, not model counts in the predictive experiment.

\artifact{src/build_matrices.py} enumerates the canonical 56 English tasks across classification, clustering, pair classification, reranking, retrieval, STS, and summarization. Only score-record language lists exactly \texttt{[eng-Latn]} or \texttt{[eng\_Latn]} qualify; cross-lingual pairs and other language lists do not. Only numeric finite \texttt{main\_score} values in $[-1,1]$ are extracted, on the designated held-out split. Eligible subset scores are macro-averaged, with no metric substitution. Candidate revisions explicitly lacking text modality or marked cross-encoder are excluded; among valid revisions the one with maximum task coverage is chosen, breaking ties by lexicographic revision name. Revisions are never stitched.

Of 290 directories with at least one usable score, 73 have fewer than 20 canonical scores and are excluded. The remaining 217 rows contain 9,832 observed cells across 56 tasks. Keeping columns observed for at least 160 rows yields the 47-task primary matrix, density 0.85185, 131 complete rows, and 86 partial historical rows. The excluded nine are AmazonCounterfactualClassification, AmazonReviewsClassification, MassiveIntentClassification, MassiveScenarioClassification, MTOPDomainClassification, MTOPIntentClassification, CQADupstackRetrieval, MSMARCO, and STS22. The exact retained task/category list is obtained by matching \artifact{data/processed/mteb_tasks.csv} to \artifact{artifacts/analysis/mteb_column_coverage.csv}. Scores on excluded tasks are not outcomes of this study. The all-56 complete cohort has only 56 rows; no optional all-56 predictive run is reported.

For OpenLLM, numeric coercion, complete raw-score filtering, range $[0,1]$, and the unflagged requirement precede exact-fullname deduplication. Within an exact fullname, priority is $2\,1\{\text{precision}\ne\text{bfloat16}\}+1\{\text{weight type}\ne\text{Original}\}$, followed by \texttt{eval\_name} for deterministic ties. Model licenses are not treated as dataset licenses. No external identity linkage or model-card crawling is used.

\subsection{Train-Only Data Flow and Grouping}
For each seed, sort the unique groups present among complete rows, shuffle with NumPy's default generator, and allocate the first $\lfloor0.6G\rfloor$ groups to training, the next through $\lfloor0.8G\rfloor$ to validation, and the rest to test, subject to minimum group boundaries in the code. Partial rows enter training only if their grouping label and canonical identity do not occur in held-out groups/identities. Native random MTEB splits contain 164/26/27 train/validation/test rows, including the 86 partial rows in training; complete-only training retains 78 rows. OpenLLM random counts are 2698/899/899 except seed 71 (2697/900/899), because canonical identity grouping controls allocation. Grouped splits can be far from 60/20/20 in row counts; exact counts appear in \artifact{artifacts/analysis/split_audit.csv}.

MTEB canonical identities lowercase metadata names and remove parenthesized suffixes, grouping embedding-dimension configurations. OpenLLM canonical identities lowercase fullnames. Publishers derive mechanically from the namespace (with the extraction fallback for MTEB names lacking one). Name-family groups scan visible names and, for OpenLLM, visible Base Model fields for ordered keyword patterns. The implemented tags include qwen, llama, mixtral, mistral, gemma, phi, deepseek, yi, falcon, olmo, smollm, internlm, solar, stablelm, e5, bge, gte, minilm, mpnet, arctic, granite, jina, and bert; unmatched names fall back to a publisher group. The first matching pattern wins. This mechanical procedure is reproducible but cannot establish verified ancestry or exclude every base/fine-tune relative.

Full-snapshot availability determines the task universe and complete-row eligibility. Only training scores determine fitted moments, imputations, factors, and regressions; validation scores tune hyperparameters and the supervised query policy. The predictor API accepts an observed matrix of exactly $|S|$ selected columns, never the full test row. Only the scoring routine reads target values. Metadata establishes split groups and provenance but is not a score-prediction feature. Manifests record columns, split identities, target identities, input hashes, and code hashes before scoring. These controls prevent direct hidden-score use in prediction but do not eliminate cohort-selection bias, exact-vector overlap, or unrecognized relatives.

The original seed-701 MTEB target set is AmazonPolarityClassification, EmotionClassification, MedrxivClusteringS2S, RedditClustering, StackExchangeClustering, SciDocsRR, DBPedia, QuoraRetrieval, SciFact, BIOSSES, STS14, and STS15. The original OpenLLM targets are IFEval and MMLU-PRO. Neither set was optimized on outcomes.

\subsection{Data QC and Descriptive Structure}
The archived test log records 13 integrity/unit tests passing in 11.23 seconds, including conditioning algebra, PSD repair, masking, split separation, selected-input prediction, and selected-target exclusion. The final artifact inspection found no missing RMSE/MAE/standardized-RMSE records, duplicate method/mode/budget/model loss keys, or fixed-target/query overlaps; query logs retain random repetitions 0--9. This is not an assertion of exhaustive software verification.

For context only, whole-snapshot complete-row Pearson task correlations have medians 0.8093 (MTEB) and 0.6860 (OpenLLM). The first three PCs of the centered, unstandardized complete matrices explain 0.9181 and 0.9810 of snapshot variance, respectively. These are post-hoc descriptive audits, not train-fitted selection inputs or guarantees of recoverability. MTEB observed-task count and snapshot standardized average quality have Spearman correlation $-0.2448$; this does not identify the missingness mechanism. The underlying audit is \artifact{artifacts/analysis/data_audit.json}.

\section{Exact Fitting, Tuning, and Masking Recipes}
\label{app:fit}
\subsection{Covariance Surrogate}
For task $j$, compute the mean $\mu_j$ and sample standard deviation on observed training entries, with scale floor $s_j=\max(\operatorname{sd}_j,0.01)$. A task with fewer than two observed training entries is nonestimable and raises an error rather than borrowing held-out data. Write $z_{ij}=(X_{ij}-\mu_j)/s_j$, $O_{ij}=1$ for observed cells, $Z_0$ for missing standardized entries replaced with zero, and $n_{ab}=(O^TO)_{ab}$. The implemented initial off-diagonal matrix is
\begin{equation}
Q_{ab}=\begin{cases}
\operatorname{clip}_{[-1,1]}\!\left((Z_0^TZ_0)_{ab}/\max(n_{ab}-1,1)\right),&n_{ab}\ge2,\\
0,&n_{ab}<2.
\end{cases}
\end{equation}
Apply $Q_{ab}\leftarrow Q_{ab}n_{ab}/(n_{ab}+\lambda)$ when $\lambda>0$ and set all diagonal entries to one. Symmetrize, eigendecompose, floor eigenvalues at $10^{-5}$, reconstruct, and divide each entry by the square root of the product of its reconstructed diagonals. Finally set $\hat R=(1-\alpha)R_{\mathrm{repair}}+\alpha I$ and $\hat C_{ab}=s_as_b\hat R_{ab}$. This training-only pairwise construction is not a pairwise-complete Pearson estimator with pair-specific centering. The diagonal-shrinkage grid keeps the primary model positive definite. Minimal PSD repair uses $\lambda=\alpha=0$ but retains the eigenvalue floor.

The nine covariance candidates are $\lambda\in\{0,10,50\}$ crossed with $\alpha\in\{0.05,0.2,0.5\}$. The no-support comparator is the best candidate on exactly the same $\alpha$ grid with $\lambda=0$. All covariance selector rules share the selected predictor hyperparameters; they are not independently test-tuned. Diagnostics in \artifact{artifacts/analysis/fit_configs.csv} include pre-repair minimum eigenvalue and negative-eigenvalue count, repair Frobenius magnitude, pair support, density, and selected parameters. On native/random MTEB the minimum pair support is 92, median support 110, and repair Frobenius magnitudes are about 1.05--1.19. Under 25\% retention median pair support is seven, repair is about 14.63--15.49, and one split contains an unobserved pair. Zero pair support contributes zero off-diagonal empirical product before repair; no evaluation score is borrowed. A zero-support \emph{pair} is distinct from a nonestimable task.

\subsection{Direct Ridge and Quality Baselines}
Direct ridge uses only selected standardized features, mean-imputing missing training predictors to standardized zero. Target $t$ is trained only on rows with observed target $t$. For $f_{iS}=(1,z_{iS})$, it solves
\begin{equation}
\min_{b_{t0},b_t}\sum_{i:O_{it}=1}(z_{it}-b_{t0}-z_{iS}^Tb_t)^2+\eta\|b_t\|_2^2,
\end{equation}
with an unpenalized affine intercept and a $10^{-10}$ numerical diagonal addition in the solver. Raw predictions undo target standardization and are clipped. No unselected or imputed hidden test features enter prediction. The ridge grid is $\eta\in\{0.01,0.1,1,10,100\}$. Validation-forward-greedy adds the candidate minimizing validation mean model RMSE after each query, either on the remaining tasks or the specified fixed target set. It reuses the same outer validation set that tuned $\eta$; no nested-CV optimality is asserted.

The quality-offset baseline averages selected standardized scores for a new model, $q_i=|S|^{-1}\sum_{j\in S}(x_{ij}-\mu_j)/s_j$, then predicts $\mu_t+s_tq_i$ and clips. The mean baseline predicts training task means while excluding queried tasks in scoring. The mean-on-SRRD-targets control uses the risk-selected query path but ignores its observed values, probing target composition rather than transferring observed information.

\subsection{Low-Rank EM-Style Imputation}
Missing standardized entries initialize to zero. For at most 25 iterations, compute an SVD of the filled matrix, form its rank-$r$ reconstruction, retain original observed entries, and update only missing entries. Stop early if the maximum absolute change is below $10^{-5}$. This is iterative low-rank reconstruction (EM-style imputation), not a guarantee of a converged maximum-likelihood EM solution. Recompute an SVD of the resulting filled matrix. Loadings are $F=V_r\operatorname{diag}(\sigma_r)/\sqrt{M_{\rm train}-1}$, and observed-entry residual variances are floored at 0.02 in standardized units. The predictor covariance is
\begin{equation}
\hat C_{\rm LR}=\operatorname{diag}(s)\,[FF^T+\gamma\operatorname{diag}(v_{\rm residual})]\,\operatorname{diag}(s).
\end{equation}
Ranks are the distinct values from $\{1,2,4,8,12\}$ capped at $n-1$ and at $\min(M_{\rm train},n)-1$ inside fitting; OpenLLM's effective candidate ranks are $\{1,2,4,5\}$. Residual multipliers are $\gamma\in\{0.1,1,5\}$. Prediction uses the same conditional formula and clipping as covariance. Random low-rank and the exploratory low-rank risk/target-risk selectors use the same tuned low-rank covariance. Many incomplete-data fits have final maximum changes well above $10^{-5}$: native/random MTEB values range 0.1038--0.2172. The fixed iteration budget is therefore an important baseline limitation; no convergence or fully optimized low-rank performance is claimed.

\subsection{Validation Objective and Random Policies}
Hyperparameters minimize validation mean per-model RMSE averaged over three deterministic random subsets at each tuning budget: $\{2,8,16\}$ for MTEB and $\{1,2,4\}$ for OpenLLM. These subsets are the first three permutations of the ten generated by \texttt{default\_rng(seed+10001)}. Selection rules do not tune their own covariance parameters. Random evaluation uses prefixes of ten full-task permutations; the fixed-target original audit filters each permutation to candidates before taking a prefix. All random-policy loss repetitions are averaged within a model before saving and aggregation.

The post-core crossed controls reconstruct the chosen predictors from archived configurations, with no expanded search: risk queries with ridge, ridge queries with covariance, and low-rank covariance design. Their historical fitting and test scores are unchanged; their exploratory timing is not concealed by presenting them as pre-outcome hypotheses.

\subsection{Training-Only Missingness Stresses}
Keep validation and test intact. Complete-only MTEB drops incomplete training rows. Random retention uses \texttt{default\_rng(seed+9001)} and independently retains each observed training entry with probability 0.5 or 0.25; already missing entries remain missing. For upper-tail deletion, sort observed values within each training column using stable sorting and remove indices from $\lfloor n_j/2\rfloor$ onward (the highest half, including the middle value for odd $n_j$). Masks are shared across methods in a condition. Each condition refits and retunes from its modified training matrix. These are controlled hypothetical deletion stresses; ``MCAR'' labels in artifact names refer to the added mask, not the natural reporting process.

\section{Uncertainty Estimand and Extended Results}
\label{app:uncertainty}
Let $I_s$ be split $s$'s test models and $\bar\ell_{sia}$ the within-model average over random query repetitions (one loss for deterministic policies). The point statistic is
\begin{equation}
\hat\theta_a=\frac15\sum_{s=1}^{5}\frac{\sum_{i\in I_s}\bar\ell_{sia}}{|I_s|}.
\end{equation}
Let $p(i)$ denote a publisher, with $P$ distinct publishers across all repeated test appearances in a dataset/split-rule/condition/mode group. For bootstrap draw $b$, sample cluster multiplicities $w^{(b)}\sim\operatorname{Multinomial}(P;1/P,\ldots,1/P)$ \emph{once jointly} across all appearances. Recompute
\begin{equation}
\hat\theta^{(b)}_a=\frac15\sum_s
\frac{\sum_{i\in I_s}w^{(b)}_{p(i)}\bar\ell_{sia}}{\sum_{i\in I_s}w^{(b)}_{p(i)}}.
\end{equation}
The percentile interval uses its 2.5th/97.5th percentiles, with 3,000 draws and seed 12345. A draw with a zero split denominator is invalid; the reported core summaries retain 3,000 valid draws. Contrasts subtract methods within the same draw. This is not an equal-identity shortcut: repeated models retain their split occurrences and each split retains equal total weight. The separate split SD is the sample SD of the five original split means. Publisher-macro summaries are supplemental diagnostics, not replacements for the primary model-weighted estimand. Publisher clusters are only a proxy for dependence; the descriptive intervals neither assume verified genealogy nor cover unmodeled refitting variation.

\subsection{Errors, Tails, and Failures}
\begin{table}[h]
\centering
\caption{Random/native diagnostic losses at the primary budgets. The tail is the equal average of five split-level 90th percentiles of per-model RMSE, not a pooled-cell quantile or a uniform error bound. All raw numbers come from \artifact{artifacts/analysis/summary.csv}.}
\label{tab:tails}
\begin{tabular}{llrrrr}
\toprule
Dataset & Method & Mean RMSE (pp) & MAE (pp) & Std. RMSE & 90th percentile (pp)\\
\midrule
MTEB ($k=8$) & SRRD & 3.4466 & 2.5505 & 0.3287 & 4.8142\\
 & Random covariance & 4.4206 & 3.1282 & 0.3503 & 5.6546\\
 & Validation-greedy ridge & 3.6253 & 2.5843 & 0.3196 & 4.9338\\
OpenLLM ($k=2$) & SRRD & 4.3552 & 3.3096 & 0.5024 & 8.0403\\
 & Random covariance & 7.1999 & 5.4051 & 0.5556 & 11.7970\\
 & Validation-greedy ridge & 4.3062 & 3.2703 & 0.4974 & 7.8677\\
\bottomrule
\end{tabular}
\end{table}
Native/random MTEB SRRD has a slightly lower raw-score mean RMSE than greedy ridge but higher training-standardized error (Table~\ref{tab:tails}); metric choice changes interpretation. Large individual failures remain. For example, the archived OpenLLM model \texttt{qingy2024/Benchmaxx-Llama-3.2-1B-Instruct} has 35.3784 pp unrun RMSE on random seed 11 and 35.4727 on publisher seed 37 at $k=2$. These are distinct appearances of one identity, not independent failures. The record does not establish whether the underlying cause is specialization, noise, or another score inconsistency. The top-100 failure file is pooled across datasets/split rules and is dominated by OpenLLM; it is not an unbiased MTEB tail sample. Full per-model losses remain in job directories, and per-task residual diagnostics for core/crossed controls are in \artifact{artifacts/predictor_controls/*_task_errors.csv}.

\begin{figure}[t]
\centering
\includegraphics[width=\textwidth]{artifacts/figures/missingness_stress.pdf}
\caption{Five-seed training-support stresses. ``Dense only'' is the 47-task complete-only historical cohort, \emph{not} a full-56-task experiment. Random retention and upper-tail deletion modify training only and trigger refitting/retuning. Native MTEB augmentation with partial rows is worse than dense-only here; support shrinkage is not uniformly beneficial. Source: \artifact{artifacts/figures/missingness_stress.pdf}.}
\label{fig:stress}
\end{figure}
\begin{table}[h]
\centering
\caption{SRRD stress results, mean RMSE $\pm$ split SD in pp. These SDs are descriptive and not standard errors.}
\begin{tabular}{lrr}
\toprule
Condition / split & MTEB ($k=8$) & OpenLLM ($k=2$)\\
\midrule
Native random & $3.4466\pm0.3298$ & $4.3552\pm0.0529$\\
Native publisher & $4.0656\pm0.3981$ & $4.3740\pm0.0628$\\
Native name-family & $3.8252\pm0.4836$ & $4.6613\pm0.3345$\\
Complete-only random & $3.1433\pm0.3047$ & Not a distinct condition\\
50\% random retention & $4.6442\pm0.6393$ & $4.3682\pm0.0639$\\
25\% random retention & $5.3492\pm0.5714$ & $4.7533\pm0.2044$\\
Upper-half deletion & $4.8816\pm0.5223$ & $6.3736\pm0.5136$\\
\bottomrule
\end{tabular}
\end{table}
The stress results reject a uniform robustness claim. In particular OpenLLM upper-tail deletion raises the split-averaged 90th-percentile RMSE to 15.6621 pp from 8.0403; MTEB raises it to 7.8454 from 4.8142. Natural missingness remains unidentified even if a particular synthetic mask is tolerated.

\subsection{Exploratory Target-Set Sensitivity}
\label{app:sensitivity}
The original seed-701 audit remains the core fixed-target result. After core outcomes, five MTEB draws with seeds 701--705 and all 15 unordered OpenLLM pairs were enumerated before additional scoring. They reuse the five frozen random/native splits and reconstruct the same fitted covariance from archived parameters without target-specific retuning. Each target set has the same candidate pool across methods and never enters it. The sensitivity random policy uses ten candidate permutations from \texttt{default\_rng(seed+10001+r)} for $r=0,\ldots,9$, rather than filtering the original stream of ten full-task permutations. Thus the deterministic seed-701 results match the core audit, but its random reference need not match exactly (e.g., MTEB 4.2597 versus 4.1875 pp). This is a separate exploratory random-policy realization, not silent replacement of the core table.

Draw losses are averaged within each model first, then summarized with exact equal-split weights. The joint publisher bootstrap uses seed 24680 with 3,000 draws. Overlapping draws/pairs and model occurrences are not independent replicates. The intervals describe the finite chosen draw average conditional on fitted predictors, splits, and target sets; they do not certify general target-population robustness.

\begin{figure}[t]
\centering
\includegraphics[width=\textwidth]{artifacts/figures/target_draw_sensitivity.pdf}
\caption{Exploratory fixed-target heterogeneity at $k=8$ (MTEB) or $k=2$ (OpenLLM). Bars show target-risk minus random covariance RMSE for each target set, sorted by contrast, not by a pre-outcome ordering. MTEB has five overlapping draws; OpenLLM enumerates every two-task pair. Source: \artifact{artifacts/figures/target_draw_sensitivity.pdf}.}
\label{fig:targetdraws}
\end{figure}
\begin{table}[h]
\centering
\caption{All five exploratory MTEB target draws at $k=8$, equal-split mean RMSE (pp). No draw was chosen to replace the original audit. Source: \artifact{artifacts/analysis/target_draws.csv}.}
\begin{tabular}{lrrr}
\toprule
Draw seed & Random covariance & Target-risk & Unrun-risk\\
\midrule
701 & 4.2597 & 3.7322 & 3.6984\\
702 & 4.8136 & 3.9746 & 4.3056\\
703 & 3.4508 & 2.8356 & 3.3670\\
704 & 4.0189 & 3.3671 & 3.8083\\
705 & 3.7406 & 3.0223 & 3.2996\\
\bottomrule
\end{tabular}
\end{table}
At $k=8$, target-risk beats the random reference on all five MTEB draws; target-risk mean errors range 2.8356--3.9746 pp. At $k=2$, target-risk beats the random reference on all 15 OpenLLM pairs but spans 2.1594--10.7364 pp. The best mean here does not imply a universal 3 pp budget. Full pair identities and values are in \artifact{artifacts/analysis/target_draws.csv}, using column IDs $0$--$5$ for IFEval, BBH, MATH Lvl 5, GPQA, MUSR, MMLU-PRO. Target-risk and unrun-risk frequently tie, especially once the full candidate pool is measured. For OpenLLM at $k=1$, their pair-averaged results differ in direction: target-risk improves on random by 1.5604 pp, whereas unrun-risk is 0.3013 pp worse. This reinforces the need to state the objective and target set rather than presenting every query policy as equivalent transfer.

\subsection{Completed Exploratory Duplicate Exclusion}
\label{app:duplicates}
The primary cohort retains the exact-vector overlaps reported in Section~\ref{sec:protocol}; identity deduplication did not eliminate them. \artifact{artifacts/analysis/split_audit.csv} counts \emph{distinct complete vectors} shared by training and test, not duplicate rows or verified aliases. Equal vectors may reflect rounding, copying, aliases, or coincidental agreement; the audit does not identify which.

A post-core sensitivity removes \emph{all} 100 OpenLLM rows belonging to repeated complete score vectors, leaving 4,396 rows with unique full vectors. The existence of \artifact{artifacts/duplicate_sensitivity/DONE.json} confirms completion across five random seeds. The existing fitting/tuning recipe is rerun on this reduced cohort, with the same budget grid and ten within-model random-policy repetitions. It uses an outcome-defined cohort filter and changes the split membership; it is not a pre-outcome clean-room validation or a verified lineage audit. At $k=2$, equal-split mean $\pm$ split SD scores are SRRD $4.3063\pm0.0638$, random covariance $7.1610\pm0.6388$, greedy ridge $4.2703\pm0.0588$, random ridge $7.1010\pm0.6325$, and low-rank design $4.2700\pm0.0586$ pp. The comparison direction persists, but this cannot retroactively make the primary results leakage-free. These are descriptive means/SDs from archived losses; no duplicate-exclusion bootstrap interval is asserted. Exact losses, excluded names, split manifests, and tuning logs are in \artifact{artifacts/duplicate_sensitivity/}; the completed run reports 6.62 seconds separately from the 65-job core accounting.

\section{Reproducibility, Resources, and Artifact Map}
\label{app:repro}
\subsection{Environment and Measured Resources}
\artifact{artifacts/environment.json} records Python 3.11.15, Linux x86\_64/glibc 2.41, 40 visible CPUs, and a four-thread numerical CPU cap. There is no GPU or paid API. The pinned \artifact{requirements.txt} lists NumPy 2.4.6, SciPy 1.17.1, pandas 3.0.6, scikit-learn 1.9.1, matplotlib 3.11.2, requests 2.34.2, pyarrow 26.0.0, pytest 9.2.0, threadpoolctl 3.7.0, and pypdf 6.21.4.

The 65 finalized core/stress jobs have a sum of recorded per-job wall times of 189.03 seconds and a peak process high-water RSS of 816,892 KiB (about 0.78 GiB). These logs include job-specific tuning/evaluation/logging, not just model fitting, and their sum is not end-to-end elapsed research time. It excludes data acquisition, the discarded-version pilots, later controls/target/duplicate checks, analysis, and manuscript compilation. Analysis itself records 93.64 seconds. High-water RSS can persist across sequential jobs in one process, so it is not an isolated allocation profile of each fit. The workspace's raw data plus outputs were reported at approximately 3.4 GB during execution, below the 5 GB budget; no precise final byte total is asserted.

\subsection{Commands and Source-to-Result Mapping}
The commands below document reproduction from the existing code. They were \emph{not} rerun to write this paper. Paths are absolute, and source scripts derive the project root from their own locations. Data acquisition requires access only to the permitted fixed data endpoints; manuscript drafting used no web search.
\begin{verbatim}
python -m pip install -r /work/requirements.txt
python /work/src/download_data.py
python /work/src/build_matrices.py
python -m pytest /work/tests
python /work/src/run_experiments.py --suite all --dataset both
python /work/src/predictor_controls.py
python /work/src/target_sensitivity.py
python /work/src/duplicate_sensitivity.py
python /work/src/theory_check.py
python /work/src/analyze.py
latexmk -pdf -interaction=nonstopmode -halt-on-error -cd /work/paper.tex
\end{verbatim}
\begin{table}[h]
\centering
\caption{Code-to-artifact map. Paths are relative to \texttt{/work}; scripts are in \texttt{src/}. Generated table/figure files are included unchanged by the manuscript.}
\begin{tabular}{p{0.30\textwidth}p{0.63\textwidth}}
\toprule
Code / source & Outputs or manuscript use\\
\midrule
\artifact{download_data.py} & \artifact{data/raw/}, \artifact{artifacts/provenance/}\\
\artifact{build_matrices.py} & \artifact{data/processed/}, cell provenance and task metadata\\
\artifact{methods.py}, \artifact{run_experiments.py} & \artifact{artifacts/results/}: manifests, tuning, losses, selections\\
\artifact{predictor_controls.py} & \artifact{artifacts/predictor_controls/}: crossed/low-rank controls, task errors\\
\artifact{target_sensitivity.py} & \artifact{artifacts/target_sensitivity/}: fixed-draw losses and selections\\
\artifact{duplicate_sensitivity.py} & \artifact{artifacts/duplicate_sensitivity/}: reduced-cohort losses/configurations\\
\artifact{theory_check.py} & \artifact{artifacts/theory/theory_mc.csv}, \artifact{artifacts/theory/theory_checks.json}\\
\artifact{analyze.py} & \artifact{artifacts/analysis/}: summaries, contrasts, thresholds and audits\\
\artifact{analyze.py} & \artifact{artifacts/tables/main_results.tex} (Table~\ref{tab:main})\\
\artifact{analyze.py} & \artifact{artifacts/tables/fixed_results.tex} (Table~\ref{tab:fixed})\\
\artifact{analyze.py} & \artifact{artifacts/tables/ablations.tex} (Table~\ref{tab:ablation})\\
\artifact{analyze.py} & \artifact{artifacts/figures/budget_curves.pdf} (Figure~\ref{fig:budget})\\
\artifact{analyze.py} & \artifact{artifacts/figures/fixed_target_curves.pdf} (Figure~\ref{fig:fixed})\\
\artifact{analyze.py} & \artifact{artifacts/figures/missingness_stress.pdf} (Figure~\ref{fig:stress})\\
\artifact{analyze.py} & \artifact{artifacts/figures/target_draw_sensitivity.pdf} (Figure~\ref{fig:targetdraws})\\
\bottomrule
\end{tabular}
\end{table}
Original full job directories are immutable/resumable through their manifests and completion markers. Analysis loads completed, current-codehash jobs with compressed losses, excluding old pilots, and asserts 65 final core configurations. Configuration hashes support provenance, not proof of absence of bugs. Sources for statements in the abstract are Tables~\ref{tab:main}--\ref{tab:ablation}, Figures~\ref{fig:budget}--\ref{fig:fixed}, and the archived paired-difference/threshold ledgers.

\subsection{Uncompleted Work and Verification Boundaries}
The optional full-56-task MTEB run, temporal/prospective tests, verified genealogy splits, natural-MNAR identification, item-level IRT reproduction, cost-aware execution, calibrated prediction intervals, and unrestricted priority review were not performed. Future extensions require separately specified protocols; no unimplemented experiment is presented as reproducible by a fictitious command.

The official ICML 2026 style and bibliography style are unmodified. The archived template check records HTTP 404 for the attempted official 2027 style URL on 2026-11-03; this artifact is not a 2027 submission. No mandatory standalone checklist was located in the inspected instructions; the checklist below is voluntary.

\section{Voluntary Reproducibility and Reporting Checklist}
\label{app:checklist}
\begin{itemize}
\item \textbf{Claims and scope: yes.} Primary embeddings and secondary generative results are separated; claims are limited to availability-selected cohorts and tested budgets (Sections~\ref{sec:protocol}--\ref{sec:discussion}). No first/SOTA or comprehensive novelty claim is made.
\item \textbf{Data and preprocessing: yes, with caveats.} Fixed commits, extraction, availability selection, excluded tasks, metric bounds, identity/group rules, and cell provenance are documented (Appendix~\ref{app:data}). Natural missingness and source contamination are not resolved.
\item \textbf{Split integrity: partial.} Predictors use only measured test coordinates and train-fitted parameters; split manifests and duplicate audits are available. Primary OpenLLM exact-vector overlap and incomplete lineage blocking remain disclosed defects/limitations, not a claim of zero leakage. The exploratory exclusion does not replace the primary cohort.
\item \textbf{Baselines and tuning: yes.} Bounded grids, selected-only ridge features, missing-data handling, validation reuse, low-rank iteration limits, and post-core exploratory controls are specified (Appendix~\ref{app:fit}). Baselines are not claimed fully optimized or nested-CV optimal.
\item \textbf{Uncertainty and multiplicity: yes, conditional/descriptive only.} Joint publisher bootstrap preserves equal-split weights and repeated appearances; split SD is separate. No independent-seed, simultaneous-coverage, confirmatory-$p$-value, or calibrated-prediction-interval claim is made (Appendix~\ref{app:uncertainty}).
\item \textbf{Negative results and failures: yes.} Unsupported support shrinkage, complete-only advantage, upper-tail degradation, substantial individual errors, fixed-target threshold failures, and uncompleted optional experiments are included.
\item \textbf{Theory and simulation: yes, under explicit assumptions.} Full derivations, row-space recovery converse, exact-arithmetic/guard/clipping qualifications, and simulated population checks are supplied (Appendix~\ref{app:proofs}). There is no proof of global greedy optimality or real-score Gaussian adequacy.
\item \textbf{Code/resources/provenance: yes, locally.} Pinned dependencies, seeds, commands, code-to-table mapping, job logs, and measured resource accounting are present (Appendix~\ref{app:repro}). A final independent end-to-end rerun and a final exact disk-size audit are not claimed.
\item \textbf{Ethics, licenses, and publication status: qualified.} Impact Statement included before references; public-source license uncertainty is disclosed. This is an anonymous research artifact, not a submitted or accepted paper.
\item \textbf{Related-work completeness: no, by imposed access constraint.} Only six permitted sources were read; no exact-problem priority claim is possible (Appendix~\ref{app:access}).
\end{itemize}

\section{Literature Access and Provenance}
\label{app:access}
Only the six references in \artifact{references.bib}, in archived v1 versions, were used; their reading notes are \artifact{artifacts/literature/allowed_literature.md}. No bibliography entry was followed, and no forbidden work or citing literature was sought or inferred. The existing access log records an earlier local command inadvertently displaying a few bibliography-line snippets from allowed PDFs. That pattern was stopped; no excluded work was identified and no snippet supplied evidence. This imposed access restriction prevents a full related-work priority check and must remain visible in later versions.

\end{document}
